\par\noindent\textbf{(1)} Hamilton量 \(H=(p_x^2+p_y^2)/(2m)+m(\omega_x^2x^2+\omega_y^2y^2)/2\)。设 \(S=W_x(x)+W_y(y)-Et\)，分离常数 \(E_x+E_y=E\)，得到 \((W_x')^2=2mE_x-m^2\omega_x^2x^2\)，另一坐标同理。分离成立的原因是势函数对两坐标可加。
\par\noindent\textbf{(2)} 定义 \(J_i=(2\pi)^{-1}\oint p_i\,dq_i\)。把相平面椭圆 \(p_i^2/(2m)+m\omega_i^2q_i^2/2=E_i\) 的面积除以 \(2\pi\)，得 \(J_i=E_i/\omega_i\)。因此 \(H(J_x,J_y)=\omega_xJ_x+\omega_yJ_y\)，作用变量在静态Hamilton系统中均为常量。
\par\noindent\textbf{(3)} 正则共轭角满足 \(\dot\vartheta_i=\partial H/\partial J_i=\omega_i\)。相空间轨道写为 \(q_i=\sqrt{2J_i/(m\omega_i)}\sin(\vartheta_i)\) 与 \(p_i=\sqrt{2m\omega_iJ_i}\cos(\vartheta_i)\)，其中 \(\vartheta_i=\omega_it+\vartheta_{i0}\)。将其代回Hamilton方程可同时验证位置、动量演化。
\par\noindent\textbf{(4)} 轨道在有限时间 \(T\) 闭合要求 \(\omega_xT=2\pi n_x,\ \omega_yT=2\pi n_y\) 对整数成立，即 \(\omega_x/\omega_y\in\mathbb Q\)；无理频率比时在相应二维相位环面上稠密。共振为退化频率关系，并不使本可分离无扰系统自身出现能量交换。
